\section{Conclusion}
\label{sec:conclusion}

The \bisect{} audit chain---three levels of SHA-256 hashes linking
Census Bureau data through computation to final plan---is a cryptographically
sound defense against all gaming vectors identified in Track R.

The security argument is simple: SHA-256 collision resistance makes
undetected manifest manipulation computationally infeasible at $2^{80}$
operations.
Any gaming attack that modifies input data, algorithm parameters, or
computation outputs changes at least one hash in the chain.
That change is detectable by any party who runs \texttt{bisect label-verify}.

The Daubert analysis confirms that this audit chain meets all four criteria
for admissible scientific evidence: testability, peer review, bounded error
rate ($2^{-80}$), and general acceptance of SHA-256.
This makes the audit chain not merely a technical safeguard but a
litigation-grade provenance instrument.

The model DIA statutory language in Section~\ref{sec:legal} translates
the technical audit chain into enforceable legal obligations: 30-day
manifest filing, 10-year retention, public verification rights, and
a burden-shifting rule for verification failures.
State legislatures can adopt this language directly.

The Track R program's conclusion follows from R.0--R.4 in combination:
when the DIA protocol is followed, the maximum achievable partisan shift
from all gaming vectors combined is $\leq 0.5$ Democratic seats nationally
(within the algorithmic noise floor), and any larger manipulation is
detectable with probability approaching 1.
Algorithmic redistricting under the DIA protocol is not merely partisan-blind;
it is verifiably partisan-blind.
